\paragraph{Failure of an abstract purity shortcut.}
Even an abelian target does not prevent a nonzero torsion section from
vanishing only in codimension two. Over
$T=\operatorname{Spec}\mathbf F_2[u,(1+u^6)^{-1}]$, let
\[
 E_u:\quad Y^2Z+u^2XYZ+YZ^2=X^3,\qquad
 O=[0:1:0],\quad P_u=[u:1:u^3].
\]
The discriminant is $1+u^6$. Directly, a singular point with $Z=1$
would satisfy $u^2X+1=u^2Y+X^2=Y^2+u^2XY=0$, forcing $u^6=1$;
the unique point at infinity is smooth. The inverse involution
$[X:Y:Z]\mapsto[X:Y+u^2X+Z:Z]$ fixes $P_u$, so $2P_u=0$,
and $P_u=O$ exactly at $u=0$. Hence over
\[
 S=\operatorname{Spec}\mathbf F_2[s,t,((1+s^6)(1+t^6))^{-1}]
\]
the homomorphism $(\mathbf Z/2)_S\to E_s\times_SE_t$ generated
by $(P_s,P_t)$ is zero precisely at the codimension-two origin.
Its target is an abelian scheme. Via the Poincare bundle this is also
a relative two-torsion Picard class on the dual abelian scheme, but
it is not a class on a family of smooth hyperplane sections of a fixed
variety. Thus it only rejects the proposed abstract group-theoretic
purity argument.

For a finite locally free target in any characteristic $p$, use
\[
 G=\operatorname{Spec}k[s,t,x,y]/
 (x^p-s^{p-1}x,\ y^p-t^{p-1}y),
\]
with $x,y$ primitive and the homomorphism $i\mapsto(is,it)$ from
$(\mathbf Z/p)_S$. Monicity gives rank $p^2$, and additivity of the
relations proves the Hopf identities. Again the zero locus is exactly
the origin. With $N$ factors the zero locus has codimension $N$.
The finite checks for $p=2,3,5,7$ and the elliptic identities are in
\texttt{check\_group\_purity.output.json}; their measured runtime is
$0.066595723$ seconds. The general claims follow from the displayed
identities, not from finite sampling.
